\documentclass[12pt]{article}
\usepackage[utf8]{inputenc}
\usepackage[T1]{fontenc}
\usepackage{hyperref}
\title{Wedau Language Resources for K-12 Curriculum Development}
\author{Compiled by Claude Sonnet 5 (claude-sonnet-5) for Andr\'as Kornai}
\date{}
\begin{document}
\maketitle
\section{Introduction}
Wedau (ISO 639-3 wed) is an Austronesian language of Milne Bay Province, Papua New Guinea. My searches found Scripture portions but no dictionary, grammar or school materials.
\section{Primary Educational Resources}
None found in search results.
\section{Linguistic and Reference Materials}
\begin{itemize}
\item Wikipedia article on Taupota, which discusses Wedau: \url{https://en.wikipedia.org/wiki/Taupota_language}
\end{itemize}
\section{Religious and Community Resources}
\begin{itemize}
\item Gospel of Mark and Acts in Wedau, with front matter on orthography (PDF): \url{https://png.bible/pdf/wed/wed_all.pdf}
\item Front matter (PDF): \url{https://png.bible/pdf/wed/wed_FRT.pdf}
\item eBible introduction page: \url{https://ebible.org/wed/INT01.htm}
\item Joshua Project profile: \url{https://joshuaproject.net/people_groups/15900}
\end{itemize}
The texts carry a Creative Commons BY-NC-ND 4.0 licence.
\section{Digital and Audio Resources}
None found in search results.
\section{Curriculum Development Context}
The Scripture texts could serve as a small corpus and as evidence for the orthography. Everything else, including a dictionary, grammar and reading materials, would need to be created.
\end{document}